% Focus: Formalize the discrete-to-continuous bridge and package the flow equation as a precise statement.

\subsection{Threadlike expressions and propagation order}\label{subsec:flow-threadlike}
% Focus: Define the class of expressions and the path language used to encode evaluation history.

We begin with arithmetic expressions over the rational numbers $\mathbb{Q}$, in order to keep the first layer of the theory combinatorial and to avoid analytic subtleties unrelated to evaluation histories.

\begin{definition}\label{def:flow-arithmetic-expression}
An \emph{arithmetic expression} over $\mathbb{Q}$ is generated by the production rules
\begin{equation}\label{eq:flow-productionrule}
\begin{aligned}
a &\longleftarrow x, \\
a &\longleftarrow (a+a), \\
a &\longleftarrow (a-a), \\
a &\longleftarrow (a\times a), \\
a &\longleftarrow (a\div a),
\end{aligned}
\end{equation}
where $x\in\mathbb{Q}$. We write $a\in \mathbb{E}[\mathbb{Q}]$.
\end{definition}

An expression may be viewed either as a formal string or as a syntax tree whose leaves are numbers and whose internal nodes are the operators $+,-,\times,\div$. Evaluation is the usual recursive process, and it is a partial function
\[
\nu:\mathbb{E}[\mathbb{Q}]\to \mathbb{Q},
\]
undefined exactly when a division by zero occurs during evaluation. Unless explicitly stated otherwise, expressions in this paper are assumed to be \emph{evaluable}, i.e.\ $\nu(a)$ is defined.

\begin{definition}\label{def:flow-evaluation-order}
The \emph{evaluation order} of an arithmetic expression is an ordering of the internal nodes in its syntax tree such that each node is evaluated only after all of its children have been evaluated.
\end{definition}

Even for elementary expressions, evaluation order need not be unique. AEG is driven by the observation that a fixed endpoint value may conceal several distinct evaluation histories.

We isolate a particularly rigid class for which evaluation histories behave like directed paths.

\begin{definition}\label{def:flow-threadlike}
A \emph{threadlike expression} is an arithmetic expression whose tree representation has the property that every left child is a leaf.
\end{definition}

Intuitively, a threadlike expression is built by starting from one initial number and repeatedly attaching a new operation and a new number on the right. Such expressions admit a canonical left-to-right reading which is convenient for geometry.

\subsubsection*{Currying and path notation}
We rewrite threadlike expressions as compositions of unary operators acting on the current value. We use the basic operators
\begin{itemize}
    \item $\oplus_y : x \mapsto x+y$,
    \item $\ominus_y : x \mapsto x-y$,
    \item $\otimes_y : x \mapsto x\cdot e^y$,
    \item $\oslash_y : x \mapsto x\cdot e^{-y}$,
\end{itemize}
where $y\in\mathbb{R}$ records the additive increment or logarithmic multiplicative increment.

For operators $a_1,a_2,\dots,a_n$, we write
\[
xa_1a_2\cdots a_n \coloneqq a_n\bigl(a_{n-1}(\cdots a_2(a_1(x))\cdots)\bigr).
\]
If a path begins with a number, we call it a \emph{bounded path}; if it consists only of operators, we call it a \emph{free path}. A bounded path evaluates to a number, while a free path evaluates to a function.

\begin{lemma}\label{lemma:flow-associative}
The path notation is associative: for paths $a,b,c$ we have $a[bc]=[ab]c$.
\end{lemma}

\begin{proof}
The statement is a direct unraveling of function composition for bounded and free paths.
\end{proof}

\begin{definition}\label{definition:flow-concatenate}
The concatenation of two paths $p_1$ and $p_2$ is their composite as functions:
\[
p_1\cdot p_2 \coloneqq p_2\circ p_1.
\]
\end{definition}

\subsubsection*{A category and a groupoid of paths}
The preceding language admits a convenient categorical packaging which will later organize equality levels and loop phenomena.

Let $X$ be a set of admissible values (for example $X=\mathbb{R}$, or a subset on which all operators under consideration are defined). Let $\mathcal{P}(X)$ be the category whose objects are elements of $X$ and whose morphisms are bounded paths: a morphism $x\to y$ is a finite sequence of operators $p$ such that $y=p(x)$. Composition is concatenation of paths, and the identity morphism at $x$ is the empty path.

When the operators allowed in paths are invertible, $\mathcal{P}(X)$ is a groupoid. In particular, the operators $\oplus_y,\ominus_y,\otimes_y,\oslash_y$ are invertible, and the inverse of a bounded path is obtained by reversing the sequence and replacing each operator by its inverse.

Free paths can be viewed as the arrow labels of an \emph{action groupoid}. Let $\mathcal{W}$ denote the set of all free paths built from the chosen operators, with multiplication given by concatenation. Each $p\in\mathcal{W}$ acts (partially, if desired) on $X$ by evaluation. The corresponding action groupoid has objects $X$ and morphisms represented by pairs $(p,x)$ with source $x$ and target $p(x)$, modulo the obvious identification of different words inducing the same function on their common domain.

This viewpoint separates two kinds of information:
\begin{itemize}
    \item a \emph{morphism} $x\to y$, which records the full evaluation history as an arrow in a path groupoid;
    \item an \emph{endpoint relation} between $x$ and $y$, obtained by forgetting which arrow was used to connect them.
\end{itemize}
Loops are endomorphisms $x\to x$ in this groupoid. They encode relations among histories and will later support the multi-layered equality program: two histories may agree as arrows, may become equal after imposing relations among loops, or may merely share the same endpoints while remaining distinct arrows.

\subsection{Additive and multiplicative infinitesimal propagations}\label{subsec:flow-propagations}
% Focus: Define generators/parameters and isolate the two basic directions with clear hypotheses.

Among threadlike expressions, the most useful bookkeeping class is the alternating one: a zigzag of multiplication and addition steps. Concretely, an alternating threadlike expression is a word of the form
\begin{equation}\label{eq:flow-alternating}
\alpha = a_1 b_1 a_2 b_2 \cdots a_l b_l,
\qquad
a_i=\otimes_{\lambda_i},\quad b_i=\oplus_{\mu_i},\quad \lambda_i,\mu_i\in\mathbb{R}.
\end{equation}
It is naturally a free path, and it becomes a bounded path once an initial value is attached.

To pass from discrete histories to a differential law, we now freeze two real parameters $\mu,\lambda\in\mathbb{R}$ and isolate the two elementary generators
\begin{itemize}
    \item $\oplus_\mu:x\mapsto x+\mu$,
    \item $\ominus_\mu:x\mapsto x-\mu$,
    \item $\otimes_\lambda:x\mapsto x\cdot e^{\lambda}$,
    \item $\oslash_\lambda:x\mapsto x\cdot e^{-\lambda}$.
\end{itemize}
These are the infinitesimal prototypes of ``add/subtract'' and ``multiply/divide''. Their defining feature is that the additive action is independent of the current value while the multiplicative action scales the current value.

The first structural invariant appears already at this discrete level: additive and multiplicative propagation do not commute. Acting on a value $x$, we have the commutator defect
\begin{equation}\label{eq:flow-torsion}
\tau(x) \coloneqq x\oplus_\mu\otimes_\lambda - x\otimes_\lambda\oplus_\mu
= \mu(e^{\lambda}-1).
\end{equation}
We call $\tau$ the \emph{arithmetic torsion}. Later sections reinterpret this commutator defect globally in the ACS and locally through contact geometry; here we record it only as a guide for what the differential limit must remember.

\subsection{Derivation of the arithmetic flow equation}\label{subsec:flow-derivation}
% Focus: Give a clean derivation; state the flow equation as a proposition; record axis-case recovery.

We now idealize the two generators as infinitesimal motions in two orthogonal directions. Let $(u,v)$ be a local orthogonal coordinate system in which the $u$-direction represents the additive generator and the $v$-direction represents the multiplicative generator. We interpret $\mu$ and $\lambda$ as the infinitesimal intensities of these two actions.

Let $s$ denote arc length along a unit-speed path and let $\theta$ be the angle that the tangent vector makes with the positive $u$-direction, so that
\[
du = \cos\theta\,ds,
\qquad
dv = \sin\theta\,ds.
\]
Let $a=a(s)$ be the current assignment value. Over an infinitesimal step $ds$, the additive generator contributes $\mu\,du$ while the multiplicative generator contributes the factor $e^{\lambda dv}$. There are two natural orderings of these two actions:
\[
a(s+ds) = (a + \mu\,du)e^{\lambda dv} + O(ds^2)
\qquad\text{or}\qquad
a(s+ds) = a e^{\lambda dv} + \mu\,du + O(ds^2).
\]
Their difference is second order, in accordance with the commutator defect \eqref{eq:flow-torsion}. Expanding $e^{\lambda dv}=1+\lambda dv+O(ds^2)$ gives
\[
a(s+ds)
= a + \mu\,du + a\lambda\,dv + O(ds^2)
= a + \bigl(\mu\cos\theta + a\lambda\sin\theta\bigr)ds + O(ds^2).
\]
Dividing by $ds$ and passing to the limit yields the fundamental differential relation.

\begin{proposition}[Arithmetic flow equation]\label{prop:flow-equation}
Along a unit-speed path whose tangent makes angle $\theta$ with the additive direction, the assignment value $a$ satisfies
\begin{equation}\label{eq:flow-equation}
\frac{da}{ds} = \mu \cos \theta + a \lambda \sin \theta.
\end{equation}
\end{proposition}

For a fixed direction $\theta$, equation~\eqref{eq:flow-equation} is a linear first-order ODE. In the axis directions $\theta\in\frac{\pi}{2}\mathbb{Z}$ it reproduces the four elementary motions:
\begin{itemize}
    \item $\theta=0$: $a(s)=a_0+\mu s$,
    \item $\theta=\frac{\pi}{2}$: $a(s)=a_0 e^{\lambda s}$,
    \item $\theta=\pi$: $a(s)=a_0-\mu s$,
    \item $\theta=\frac{3\pi}{2}$: $a(s)=a_0 e^{-\lambda s}$.
\end{itemize}
In this sense, the flow equation is the infinitesimal envelope of the discrete generator calculus encoded by $\oplus_\mu$ and $\otimes_\lambda$.

\subsection{Lie-theoretic formulation via \texorpdfstring{$\mathrm{Aff}(1)$}{Aff(1)}}\label{subsec:flow-aff1}
The flow equation admits a convenient Lie-theoretic encoding which separates the group-level propagation from its infinitesimal generator. We use the following notational convention throughout: a Roman letter such as $G$ denotes a Lie group, while the corresponding fraktur letter such as $\mathfrak{g}$ denotes its Lie algebra.

\subsubsection*{The affine group and its Lie algebra}
Let
\[
G \coloneqq \mathrm{Aff}(1)
=\left\{
\begin{pmatrix}
\Phi & a\\
0 & 1
\end{pmatrix}
\;\middle|\; \Phi\in\mathbb{R}_{>0},\ a\in\mathbb{R}
\right\},
\]
acting on $x\in\mathbb{R}$ by $x\mapsto \Phi x + a$. The group law is matrix multiplication, equivalently
\[
(\Phi_1,a_1)(\Phi_2,a_2)=(\Phi_1\Phi_2,\ a_1+\Phi_1 a_2).
\]
This semidirect-product rule already mirrors propagation: a multiplicative scaling $\Phi_1$ amplifies a subsequent additive translation $a_2$.

Its Lie algebra is
\[
\mathfrak{g}\coloneqq \mathfrak{aff}(1)
=\left\{
\begin{pmatrix}
\alpha & \beta\\
0 & 0
\end{pmatrix}
\;\middle|\; \alpha,\beta\in\mathbb{R}
\right\},
\]
with bracket given by the matrix commutator. A convenient basis is
\[
H\coloneqq
\begin{pmatrix}
1 & 0\\
0 & 0
\end{pmatrix},
\qquad
N\coloneqq
\begin{pmatrix}
0 & 1\\
0 & 0
\end{pmatrix},
\qquad
[H,N]=N.
\]
The one-parameter subgroups $\exp(tH)=\begin{psmallmatrix}e^t&0\\0&1\end{psmallmatrix}$ and $\exp(tN)=\begin{psmallmatrix}1&t\\0&1\end{psmallmatrix}$ encode, respectively, pure multiplication and pure addition. In particular,
\[
\rho(\otimes_\lambda)=\exp(\lambda H),
\qquad
\rho(\oplus_\mu)=\exp(\mu N),
\]
matching the discrete generators used in \S\ref{subsec:flow-propagations}.

More generally, any free path $p$ built from these generators determines (when it is defined) an affine transformation of the line, hence an element $\rho(p)\in G$. The resulting map $\rho$ is a concrete realization of the ``word $\to$ operator'' passage implicit in the path notation. The differential viewpoint corresponds to passing from short words to their infinitesimal generators: for example, $\exp(\varepsilon H)$ and $\exp(\varepsilon N)$ encode the elementary multiplicative and additive updates to first order in $\varepsilon$, and the commutator defect recorded in \eqref{eq:flow-torsion} appears as a second-order obstruction to commutativity.

\subsubsection*{A left-invariant evolution and the flow equation}
Fix a direction $\theta$. We prescribe the Lie-algebra element
\begin{equation}\label{eq:flow-omega}
\Omega(\theta)\coloneqq \lambda\sin\theta\,H+\mu\cos\theta\,N
=
\begin{pmatrix}
\lambda\sin\theta & \mu\cos\theta\\
0 & 0
\end{pmatrix}.
\end{equation}
Consider the left-invariant linear system on $G$,
\begin{equation}\label{eq:flow-aff-ode}
\frac{d}{ds}g(s)=\Omega(\theta)\,g(s),
\qquad g(s)\in G.
\end{equation}
Writing
\[
g(s)=
\begin{pmatrix}
\Phi(s) & a(s)\\
0 & 1
\end{pmatrix},
\]
equation~\eqref{eq:flow-aff-ode} becomes
\[
\begin{pmatrix}
\Phi'(s) & a'(s)\\
0 & 0
\end{pmatrix}
=
\begin{pmatrix}
\lambda\sin\theta & \mu\cos\theta\\
0 & 0
\end{pmatrix}
\begin{pmatrix}
\Phi(s) & a(s)\\
0 & 1
\end{pmatrix}.
\]
Comparing entries yields
\[
\Phi'(s)=\lambda\sin\theta\,\Phi(s),
\qquad
a'(s)=\mu\cos\theta+a(s)\lambda\sin\theta.
\]
The second equation is exactly Proposition~\ref{prop:flow-equation}. Thus the assignment component of the arithmetic flow is the translation component of a left-invariant evolution on the affine group.
